\chapter{Limitations}
\label{ch:limitations}

\section{The Population}
\label{sec:lim-population}

\label{sec:lim-finscope}%
No \ac{FinScope} wave was available for 2017, so the 2019 wave supplies the four
financial-inclusion flags for a population otherwise fixed at 2017. The imported variables are
categorical and need no deflation, but financial inclusion was changing over that period and banked
status gates \ac{BNPL} eligibility in Submodel~12. Correcting for the gap would require an
assumption about the direction and pace of that change which this thesis cannot source.

\label{sec:lim-savings}%
\texttt{liquid\_savings} is the least reliable quantity carried into the model, deriving from a
\ac{NIDS} field that measures reported financial assets rather than an accessible cash balance, and
requiring winsorisation to remove implausible outliers. Validation pattern~4 measures the share of
agents reaching zero liquid savings, so the model is judged partly on its handling of its own
weakest input.

\label{sec:lim-servicing}%
Servicing is constructed rather than observed (Section~\ref{sec:servicing}), on three assumptions.
Statutory maximum prescribed rates are ceilings and the \ac{CCMR} publishes no realised rates, so
opening servicing is an upper bound; loan terms are assumed, with a six-month floor; and the
product mix implied by assigning sub-sector rates to a consolidated balance is itself an
assumption. The 53 households whose income falls below the Regulation 23A expense norm were
retained, and the data cannot establish why they hold the debt they report.

\label{sec:lim-reg23a}%
Regulation 23A prescribes a minimum-expense norm for an individual consumer, and the model applies
it to household income because \ac{NIDS} is a household survey. For a multi-earner household this
makes the test more permissive than the regulation would be in practice, so the model's lender is
somewhat more willing to lend than a compliant lender would be.

\label{sec:lim-resample}%
Weighted resampling raises the per-capita income Gini from 0.651 to 0.671, passing its tolerance
with little margin and, if anything, modestly overstating distress among the lowest-income agents.
The distortion is measured and not corrected.

\section{The Validation Strategy}
\label{sec:lim-validation}

\label{sec:lim-calibration}%
The most important limitation in the thesis is that the baseline is calibrated rather than
validated. The income-shock probability of Submodel~1 cannot be estimated from the data, since
\ac{NIDS} Wave 5 enters this design as a single wave and within-household income volatility is not
identified from one wave, so it is fitted until baseline arrears reproduce pattern~2. Agreement
with the 2017 \ac{CCMR} is consequently not evidence that the model is right, and only the results
obtained after \ac{BNPL} is injected, with that parameter held fixed, are genuine out-of-sample
predictions.

\label{sec:lim-units}%
Pattern~2 compares the model's arrears profile against \ac{CCMR} figures computed on an account
basis, while the model's unit is the household. An account-level arrears share is not a
household-level distress rate, so the comparison asks only whether the profile is of the right
shape and scale, and because pattern~2 is also the calibration target that imprecision propagates
into the fitted shock probability.

\label{sec:lim-targets}%
Only pattern~2 is both South African and vintage-matched to the 2017 population. Pattern~1 rests on
United States banking data, pattern~3 on a United Kingdom credit-card model, and pattern~4 on a
South African survey fielded in 2025, which is the cross-period comparison that
Section~\ref{sec:injection} otherwise forbids. Pattern~4 also compares a realised stock condition
against a stated expectation, with no established mapping between the two, so it is best read as a
plausibility check.

\label{sec:lim-pattern1}%
Pattern~1 is not fully independent of the design. The want-driven borrowing path of Submodel~3 was
adopted partly because a purely shortfall-driven \ac{BNPL} agent would treat the product as a pure
substitute and could not reproduce the deHaan et al. complementarity result under any
parameterisation. The direction of the effect is therefore partly built into the lender-choice
rule. Its magnitude is not, and the pattern can still fail if the Regulation 23A gate crowds out
traditional borrowing sufficiently.

\section{The Model}
\label{sec:lim-model}

\label{sec:lim-transfer}%
The behavioural micro-foundations of Section~\ref{sec:lit-behavioural} are drawn from United States
evidence and applied to South African borrowers: the minimum-payer share of 0.29, the present-bias
results, and the Ackert et al. experiment that motivates the want-driven borrowing path. No
developing-market evidence of comparable quality was located for any of the three. The transferred
parameter that matters most is swept from 0.20 to 0.40, and the mechanisms transferred are claims
about general features of decision-making under credit contracts rather than about United States
institutions, but neither argument establishes that the parameters are right for South Africa.

\label{sec:lim-uncalibrated}%
The peer-influence coefficient $\beta$ and the spontaneous adoption rate $q_{\text{base}}$ of
Equation~\ref{eq:peer} are fixed by no South African source. The $\beta = 0$ control arm addresses
the circularity this creates but not the plausibility problem, since it does not identify which
region of the $(q_{\text{base}}, \beta)$ surface corresponds to the actual South African market.
The available anchors are indirect: the 83.1\% banked ceiling, the TransUnion finding that 20\% of
consumers intending to apply for credit named \ac{BNPL}, and the \ac{CFPB} stacking shares.

\label{sec:lim-amount}%
Submodel~4 has no literature anchor, the sweep having located no work that fixes a
borrowing-amount rule. Sensitivity analysis across three specifications is mandatory for this rule,
and any material movement across them is reported as amount-rule sensitivity.

\label{sec:lim-structural}%
Six exclusions were adopted by design, and each bounds what the model can be asked.

\begin{itemize}
    \item \textbf{No macroeconomic dynamics:} Submodel~16 holds the environment static, since any
    time-variation would confound the injection effect with a macro effect, so the model cannot
    speak to how \ac{BNPL} stress interacts with a downturn.
    \item \textbf{No contagion:} one aggregate lender with no balance sheet that can fail, and no
    feedback from household default to income. The only household-to-household channel transmits
    \ac{BNPL} adoption, so the model produces correlated individual defaults and not cascades.
    \item \textbf{No competition among traditional lenders:} the traditional side is a single
    non-adaptive stub, suppressing the lender-conduct mechanism Hamill et al. show to matter.
    \item \textbf{No learning and no scarring:} rules are not updated from experience, and default
    carries no rehabilitation path beyond cut-off.
    \item \textbf{No peer effects on consumption:} Cardaci's expenditure-cascade mechanism is not
    implemented, and peer influence acts only on \ac{BNPL} adoption.
    \item \textbf{Static household composition:} households do not form, dissolve or change size,
    and reference groups are fixed at initialisation.
\end{itemize}

\section{How Results Should Be Read}
\label{sec:lim-reading}

Results are directional and mechanistic. They describe what happens to a population of this kind
when a lender class with these properties is introduced, and are informative about the sign, the
rough magnitude and the mechanism of that change. They are not forecasts, they do not identify a
threshold value of \ac{BNPL} access that applies outside the model, and no number in
Chapter~\ref{ch:results} should enter a policy calculation without the qualifications above. What
survives all of them is the comparison the model was built to make, between two runs identical
except for the presence of a bureau-reporting obligation, which holds every uncalibrated parameter
fixed across both arms.
